\subsection{Contrato del reto y unidad estadística}

La primera decisión técnica fue separar con precisión la \emph{unidad de observación de las
fuentes} de la \emph{unidad supervisada de predicción}. La unidad supervisada es la parcela,
identificada de forma única por \artifact{ID_POLIGONO}. El archivo oficial de
rendimiento/split contiene 197 parcelas: 138 filas de \artifact{ENTRENAMIENTO} con
\artifact{RENDIMIENTO_T_HA} observado y 59 filas de \artifact{PREDICCION} con el objetivo
reservado por FIRA.

\begin{table}[H]
\centering
\caption{Contrato básico del problema supervisado.}
\label{tab:cp01-contract}
\begin{tabular}{lr}
\toprule
Elemento & Valor \\
\midrule
Parcelas totales & 197 \\
Parcelas etiquetadas & 138 \\
Parcelas objetivo & 59 \\
Estados & Hidalgo, Puebla, Tlaxcala \\
Identificador canónico & \artifact{ID_POLIGONO} \\
Variable objetivo & \artifact{RENDIMIENTO_T_HA} \\
Ciclo objetivo & abril--octubre de 2025 \\
\bottomrule
\end{tabular}
\end{table}

El hecho de que BASIC y PRO contengan decenas de miles de filas no aumenta el tamaño muestral
supervisado. Una fila BASIC corresponde conceptualmente a una observación
parcela--fecha--sensor; una fila PRO tiene la misma naturaleza longitudinal. Por tanto,
\[
n_{\mathrm{supervisado}}=138,
\]
no $107{,}666$ ni $47{,}804$. Tratar las filas temporales como observaciones independientes
con la etiqueta parcelaria replicada produciría pseudorreplicación severa, correlación
intra-parcela ignorada y leakage entre train y validación. Toda la arquitectura de datos
posterior conserva esta distinción.

\subsection{Objetivo estadístico inicial y evolución posterior}

Durante los Checkpoints~01--03 se mantuvo una formulación convencional de regresión:
construir una representación parcelaria $x_i$ y estudiar modelos
\[
\hat y_i=f_\theta(x_i)
\]
bajo validación espacialmente consciente. Esta etapa era necesaria para establecer qué
fuentes, representaciones y familias de modelos eran razonables antes de explotar la
estructura fija del conjunto objetivo.

Posteriormente, en el Checkpoint~04, la doctrina se refinó a una formulación
transductiva. Esa reformulación no invalida los checkpoints previos: los convierte en
evidencia de ingeniería, representación y baseline. La diferencia central es que en el
problema final conocemos de antemano los 59 vectores $X_{\Target}$ que serán evaluados. Por
ello, el objeto final no es una función universal para parcelas futuras sino el vector
\[
\widehat{\bm y}_{\Target}\in\R^{59}.
\]

\subsection{Frontera de información}

Desde el primer checkpoint se mantuvo una separación estricta entre:
\begin{enumerate}
\item información observable y permitida;
\item información derivada de etiquetas visibles;
\item etiquetas reservadas.
\end{enumerate}

Para los 59 objetivos:
\[
X_{\Target}\;\text{es observable},\qquad
\bm y_{\Target}\;\text{es no observable}.
\]
La regla se aplica no sólo al fit final sino también a selección de variables,
hiperparámetros, elección de vecinos, thresholds de routing y decisiones de ensamble.

En los checkpoints convencionales, transformaciones entrenables como imputación, escalado,
PCA, PLS y descubrimiento supervisado se ajustan dentro de cada fold. En la fase transductiva
se permite que transformaciones estrictamente $X$-only conozcan $X_{\Target}$, pero nunca
$\bm y_{\Target}$. Esta distinción se vuelve especialmente importante para grafos,
estandarización transductiva, densidad y selección de pseudo-competiciones.

\subsection{Horizonte de información: clean frente a competition}

La ingeniería de datos reconoció desde el principio que una variable puede ser válida para la
competencia y, sin embargo, no corresponder a un sistema de pronóstico prospectivo. Por eso se
crearon dos modos conceptuales.

\paragraph{clean.}
Incluye información histórica, estática o disponible bajo una interpretación prospectiva
conservadora. Su objetivo es permitir una lectura científica de generalización sin usar
resultados contemporáneos que sólo aparecen después del ciclo.

\paragraph{competition.}
Incluye el conjunto clean más información pública/transductiva permitida para reconstruir el
rendimiento de las 59 parcelas después de observar el ciclo 2025. Ejemplos: series 2025
completas, WaPOR 2025, clima 2025 y, posteriormente, SIAP 2025 como evidencia municipal
contemporánea.

La separación no equivale a \emph{seguro} frente a \emph{leakage}. Una variable competition es
válida si procede de una fuente pública permitida y no codifica los 59 rendimientos ocultos.
La categoría existe para conservar honestidad semántica sobre el momento en el que la
información estaría disponible en un escenario de despliegue distinto.

\subsection{Implicación small-\texorpdfstring{$n$}{n}/high-\texorpdfstring{$p$}{p}}

Desde el inicio se anticipó que la tabla final sería de dimensión elevada respecto al número
de etiquetas. Esta condición condiciona todo el proyecto:
\[
n_{\Train}=138,\qquad p\gg n.
\]
Consecuencias:
\begin{itemize}
\item cualquier selección guiada por $y$ tiene alta varianza;
\item un grid de hiperparámetros muy grande puede sobreoptimizar los folds;
\item los errores de validación deben leerse junto con dispersión y protocolos espaciales;
\item Ridge y PLS son baselines fundamentales, no modelos demasiado simples;
\item árboles/boosting requieren fuerte disciplina de validación;
\item una mejora de milésimas en el mismo conjunto de desarrollo no debe interpretarse como
ganancia real sin un control split-excluded.
\end{itemize}

Esta última regla resultó decisiva en Checkpoint~04F: un blend Local/Graph obtuvo un mínimo
ligeramente inferior al LocalRidge en la tabla completa de desarrollo, pero no se promovió
porque la selección leave-one-split-out no sostuvo la ventaja.

\subsection{Definición operacional de éxito}

El proyecto no busca maximizar una métrica abstracta de representación. Cada checkpoint debe
reducir incertidumbre sobre una decisión concreta. En las primeras fases, la pregunta es si
una fuente puede integrarse sin ambigüedad y si una representación mantiene señal bajo
validación. En la fase transductiva, el criterio se hace más estricto:

\begin{quote}
¿La información observable de las 197 parcelas permite reducir el error en pseudo-competiciones
que reproducen la geometría de las 59 parcelas reales, sin acceder en ningún momento a sus
etiquetas ocultas?
\end{quote}

Esta formulación conecta ingeniería de datos, estadística espacial, aprendizaje supervisado y
semi-supervisado en un único contrato de decisión.
